\chapterimage{figs/chapterhead.png}
\chapter{Known Limitations}
\label{chap:appendix-known-limitations}

This appendix records the main limitations that remain explicit in the current
manual. It is not a complaint list; it is a guardrail against overclaiming.
\index{limitations}

\section{Simulation and numerical limits}

The manual documents discrete-event simulation, continuous solvers, Petri
nets, cellular automata, and whole-cell extensions, but it does not pretend
that all of those domains have the same level of validation. The discrete-event
and simulation background is grounded in standard texts \cite{banks2005,law2015};
the numerical ODE background is grounded in standard numerical methods
\cite{hairer1993}; and the Petri-net and cellular-automaton surfaces are
presented as modeling formalisms, not as blanket scientific proof of a given
result \cite{murata1989,wolfram2002}.

\section{Biological and domain limits}

The biological extension surface is documented conservatively. Flux balance
analysis is a modeling technique for metabolite flow, not a universal proof of
biological correctness \cite{orth2010}. SBML is a model exchange format, not a
validation oracle \cite{hucka2003}. Whole-cell modeling remains a demanding
research area with a specific scientific scope \cite{karr2012}. The manual
keeps that boundary visible on purpose.

\section{Build and toolchain limits}

The repository uses CMake and Qt as its current build and application
frameworks. Their official documentation is the reference point for build and
GUI semantics \cite{cmake-docs,qt-docs}. The manual does not claim that every
target or preset is equally mature, and it does not promote a build-system
example into a supported product unless another chapter says so explicitly.

\section{Testing and quality limits}

Software testing has its own evidence level. General software engineering and
testing references are useful background, but they do not turn one successful
run into a release gate \cite{sommerville2015,myers2011}. The current manual
keeps that distinction visible because the repository still separates build,
startup, functional, numerical, scientific, and release evidence.

\section{What the appendix does not claim}

\begin{itemize}
\item It does not declare public release readiness.
\item It does not declare that every experimental model is scientifically
  validated.
\item It does not declare that loopback or local-only services are safe for
  public deployment.
\item It does not declare that every current target has equal maturity.
\item It does not replace the chapter-specific validation evidence.
\end{itemize}

The manual should be read as a current-state reflection of the repository, not
as a permanent promise about future support.
